% This file should be replaced with your file with thesis content.
%=========================================================================
% Authors: Michal Bidlo, Bohuslav Křena, Jaroslav Dytrych, Petr Veigend and Adam Herout 2019

% For compilation piecewise (see projekt.tex), it is necessary to uncomment it and change
% \documentclass[../projekt.tex]{subfiles}
% \begin{document}

\chapter{Introduction}

TODO

\chapter{Abstract}
\label{abstrakt}

TODO 

\chapter{Theoretical Foundations}
\label{chap:theoretical-foundations}

\section{Prerequisites}

This thesis assumes familiarity with fundamental concepts in modern deep learning, in particular convolutional 
neural networks (CNNs), gradient-based optimisation, and automatic differentiation as implemented in frameworks 
such as PyTorch~\cite{PaszkePyTorch, PyTorchDocs}.  
Comprehensive introductions to these topics are provided in Goodfellow et al.~\cite{GoodfellowDeepLearning} and in 
the official PyTorch documentation~\cite{PyTorchDocs}.  

The remainder of the thesis is self-contained with respect to cellular automata (CA) and their neural extensions.

\section{Cellular Automata}

Cellular automata (CA) are discrete dynamical systems composed of a regular grid of cells, each of which occupies 
one state from a finite set. The grid evolves in parallel over discrete time steps according 
to a local update rule that depends solely on the states of neighbouring cells~\cite{MarinescuComplexSystems}.  

Formally, a cellular automaton is defined by:
\begin{itemize}
    \item a \textbf{lattice} of cells, typically arranged in one, two, or three dimensions,
    \item a \textbf{finite state set} $\mathcal{S}$,
    \item a \textbf{neighbourhood function} $\mathcal{N}$ describing which cells influence a given cell,
    \item an \textbf{update rule} $f : \mathcal{S}^{|\mathcal{N}|} \rightarrow \mathcal{S}$ that determines the next state~\cite{IlachinskiCellular}.
\end{itemize}

The foundational work of Ulam and von Neumann in the 1940s established CA as a model for self-replication~\cite{VonNeumann1966}.  
Later, Wolfram proposed a widely used behavioural classification of CA into four qualitative classes~\cite{Wolfram1983, Wolfram2002}. 
Class 1 systems quickly converge to a uniform state. Class 2 systems evolve into static or periodically oscillating patterns. 
Class 3 systems exhibit chaotic and seemingly unpredictable behaviour. Class 4 systems demonstrate complex dynamics characterised 
by long-lived structures capable of interaction and computation.

These classes remain a reference point in subsequent work on both classical and neural CA.

\subsection{One-Dimensional Cellular Automata}

The simplest non-trivial CA operate on a one-dimensional array of cells. In the elementary CA family, each cell observes 
its own state and that of its two immediate neighbours (radius 1, neighbourhood size 3). With binary states, the update rule 
has $2^{3} = 8$ possible neighbourhood configurations, yielding $2^{8} = 256$ distinct rules in Wolfram's enumeration scheme~\cite{Wolfram1983}.

A rule number encodes the next-state lookup table in binary. Notable examples include Rule 90, which generates 
the well-known Sierpiński triangle from a single active seed (Class 3), and Rule 30, whose chaotic evolution 
has been used as a source of pseudorandomness. Rule 110 (Class 4) is the smallest known Turing-complete 1D CA, 
supporting gliders and glider collisions that can perform universal computation~\cite{Cook2004}.

\subsection{Two-Dimensional Cellular Automata}

In 2D, the Moore neighbourhood (a $3 \times 3$ block) is the most common. Because the space of lookup tables is 
large ($2^9$ input combinations), rules are often expressed using birth/survival notation.

Conway's Game of Life~\cite{Gardner1970} remains the most influential 2D CA, featuring a wide variety of emergent structures, 
including still lifes (e.g. block, beehive), oscillators (e.g. blinker, toad), spaceships (e.g. glider, lightweight spaceship), 
and more complex constructions such as glider guns and other computational structures.

Life is Turing-complete~\cite{Berlekamp1982} and has influenced much of the later work on self-organising neural CA models.

\subsection{Three-Dimensional Cellular Automata}

In three dimensions, the Moore neighbourhood includes up to 26 adjacent cells (face-, edge-, and corner-adjacent) in a 3$\times$3$\times$3 cube.  
With binary states, this creates $2^{27} \approx 134$ million possible neighbourhood configurations, so practical 3D CA typically use 
totalistic or outer-totalistic rules that depend only on the sum of live neighbours (not their exact positions)~\cite{Wolfram2002}.

Notable examples include von Neumann's 29-state self-replicating automaton~\cite{VonNeumann1966}, 
excitable-media rules such as Brian's Brain extended to 3D~\cite{DeKinderen1998}, 
and three-dimensional Game of Life variants capable of producing branching and crystal-like growth patterns.

Advances in computing hardware in the 2000s facilitated larger-scale 3D CA studies~\cite{Tobias2009}.  
These volumetric models later aligned naturally with neural CA research.

\subsection{Differentiability and Continuous Cellular Automata}

Although this thesis uses a discrete 3D lattice, the update rule is fully differentiable, enabling gradient-based optimisation~\cite{Bena2025}.  
This connects neural CA to continuous cellular automata, where states are real-valued and update rules are smooth rather than discrete.

\textbf{SmoothLife}~\cite{Rafler2011} generalises the Game of Life to continuous space and time using continuous averaging kernels and 
smooth transition functions, producing gliders and complex wave-like structures on a continuous field.

\textbf{Lenia}~\cite{Chan2019} further extends this idea using parameterised convolution kernels and growth functions.  
Its exploration of parameter space led to hundreds of stable, self-organising patterns resembling an artificial ecology.

These systems demonstrate that complex self-organisation does not depend on discrete state transitions. Instead, strictly local, differentiable 
interactions are sufficient to generate stable and adaptive spatial structures. Neural cellular automata leverage this property while introducing 
optimisation through gradient descent, allowing update rules to be learned for specific structural objectives rather than discovered through manual 
parameter exploration.

\section{Neural Cellular Automata}

Neural Cellular Automata (NCA) were introduced by Mordvintsev et al.~\cite{Mordvintsev2020} as a differentiable 
generalisation of classical cellular automata. In their formulation, the discrete lookup-table update rule of traditional 
CA is replaced by a small neural network whose parameters are shared across all cells. While the spatial lattice remains 
discrete, the state of each cell becomes a real-valued vector, enabling gradient-based optimisation through time.

Let the lattice be denoted by $\Omega \subset \mathbb{Z}^d$, where $d$ represents the spatial dimension. 
Each cell $x \in \Omega$ maintains a state vector
\[
\mathbf{s}_t(x) \in \mathbb{R}^{C},
\]
where $C$ is the number of channels. 

The global configuration at time $t$ can therefore be viewed as a tensor containing the states of all cells.

Following Mordvintsev et al.~\cite{Mordvintsev2020}, the update process consists of two steps. 
First, a fixed perception operator extracts local neighbourhood information using convolutional filters. 
Second, a shared neural network processes this local representation and produces a state increment that is added to the current state.

This formulation preserves the defining property of cellular automata - strict locality - while replacing the hand-crafted transition 
function with a learnable parametric model. Because the update rule is differentiable with respect to its parameters, the automaton can be 
trained by unrolling its dynamics over multiple steps and applying gradient descent to minimise a loss defined on the resulting configuration.

Neural cellular automata can therefore be understood as spatially distributed dynamical systems whose behaviour is shaped by optimisation rather 
than by manually specified rules.

\section{Extension from Two to Three Dimensions}

The original neural cellular automaton proposed by Mordvintsev et al.~\cite{Mordvintsev2020} operates on a two-dimensional grid. 
Extending this framework to three-dimensional voxel lattices is conceptually straightforward but introduces several 
practical and dynamical considerations.

In two dimensions, perception is typically performed using $3 \times 3$ convolutional filters applied to a Moore neighbourhood. 
In three dimensions, this naturally generalises to $3 \times 3 \times 3$ volumetric filters, increasing the neighbourhood size 
from nine to twenty-seven cells. As demonstrated in prior 3D extensions~\cite{Sudhakaran2021,Rossi2021,Aadityaza2024}, Sobel-type 
gradient filters and other perception operators can be extended along the $x$, $y$, and $z$ axes without altering the overall update formulation.

However, volumetric grids introduce important differences. The number of cells grows cubically with lattice resolution, which substantially increases 
memory requirements and computational cost. Moreover, since updates remain strictly local, information must propagate only through repeated applications of 
the update rule. In larger three-dimensional grids, maintaining a coherent global structure therefore requires stable dynamics over a longer period of time.

These considerations suggest that behaviours observed in two-dimensional neural cellular automata - particularly persistence and regeneration - may 
not manifest identically in volumetric domains. While similar mechanisms can, in principle, be extended to three dimensions, the change in spatial 
topology and neighborhood structure can lead to differences in stability, information propagation, and emergent dynamics.

\section{Core Components of a 3D Neural Cellular Automaton}
\label{sec:core-components}

The preceding sections introduced classical cellular automata, continuous extensions, and neural CA models. 
This section focuses on the principal theoretical components that form the foundation of a three-dimensional 
neural cellular automaton. The aim is to clarify the general concepts that govern these systems without 
delving into implementation-specific details or experimental results.

\subsection{State Representation}

In neural cellular automata, each voxel maintains a vector of real-valued channels. This notion originates 
from the two-dimensional model of Mordvintsev et al.~\cite{Mordvintsev2020}, where multi-channel states 
support both visible pattern information and internal memory.

Channels can generally be divided into two categories:
\begin{itemize}
    \item \textbf{Visible channels}, representing quantities directly compared to the target pattern, such as 
    colour, occupancy, or density;
    \item \textbf{Hidden channels}, providing additional internal memory and facilitating information 
    propagation between neighbouring cells.
\end{itemize}

If each cell contains $C$ channels, its state can be represented as
\[
    \mathbf{s} \in \mathbb{R}^{C}.
\]

Many neural CA models also include a continuous \emph{alive} channel, which acts as a gating mechanism 
modulating whether updates are applied to a cell. This concept appears in the original NCA formulation 
\cite{Mordvintsev2020} and has conceptual parallels in continuous CA such as SmoothLife~\cite{Rafler2011} 
and Lenia~\cite{Chan2019}, where real-valued states enable smooth evolution and gradient-based optimisation.

\subsection{Perception Mechanism}

Despite employing continuous state vectors, neural CA retain the defining property of locality: each cell 
updates based solely on information from its immediate neighbourhood~\cite{IlachinskiCellular, Wolfram2002}. 
In neural CA, locality is realised by applying a fixed set of convolutional filters (e.g., $3 \times 3$ in 2D or 
$3 \times 3 \times 3$ in 3D) to extract spatial features~\cite{Mordvintsev2020, Sudhakaran2021}.

Typical filter types include:
\begin{itemize}
    \item \textbf{Identity filters}, capturing a cell's own state;
    \item \textbf{Neighbour aggregation filters}, such as sums or averages of adjacent voxels;
    \item \textbf{Gradient filters}, including Sobel-type kernels extended to 3D, approximating directional differences 
    and providing a local analogue of biological gradient sensing~\cite{Mordvintsev2020, Sudhakaran2021}.
\end{itemize}

The outputs of all applied filters are concatenated into a single perception vector for each cell, which is then 
passed to the update function. Using fixed, non-learnable filters stabilises training and maintains conceptual 
continuity between neural and classical cellular automata~\cite{Mordvintsev2020, Chan2019}.

\subsection{Update Rule}

The update rule is a shared, parametric function that governs how each cell's state evolves over time. Following 
the formulation of Mordvintsev et al.~\cite{Mordvintsev2020}, the rule operates on the perceived neighbourhood 
features and produces a state increment added to the current state:
\[
\mathbf{s}_{t+1} = \mathbf{s}_t + \Delta \mathbf{s}.
\]

This approach preserves the parallel, local nature of classical cellular automata while allowing the transition 
function to be learned through gradient-based optimisation. The update function is typically implemented as a 
lightweight neural network applied independently to each cell, such as a small multilayer perceptron or a sequence 
of $1 \times 1$ (or $1 \times 1 \times 1$ in 3D) convolutions~\cite{Mordvintsev2020}. These designs ensure that 
updates depend only on locally perceived information while keeping the parameter count sufficiently low to maintain 
stable training.

Two mechanisms commonly appear in neural CA formulations. Alive gating regulates updates based on the continuous 
alive channel, providing a soft analogue of classical alive/dead states and improving long-term stability~\cite{Mordvintsev2020}. 
Stochastic update masking applies updates probabilistically, preventing rigid synchronisation and encouraging 
robustness~\cite{Mordvintsev2020}.

These components collectively support behaviours such as growth, pattern persistence, and regeneration, phenomena 
also observed in continuous morphogenetic systems such as Lenia~\cite{Chan2019}. Early three-dimensional neural CA 
models~\cite{Rossi2021,Sudhakaran2021,Aadityaza2024} demonstrate that the underlying update-rule principles generalise 
naturally to volumetric grids, although the emergent dynamics may differ from their 2D counterparts.

\subsection{Grid Management}

The neural CA is defined on a finite three-dimensional lattice, with boundary conditions chosen according to the 
experiment~\cite{Mordvintsev2020}. Key considerations include the lattice dimensions, initialization, and boundary treatment.
Periodic or zero-value padding at the edges is often used to reduce boundary artefacts, though specific lattice sizes and boundary 
choices vary across implementations~\cite{Niklasson2021}.

Three-dimensional grids scale cubically with lattice size, creating computational challenges. Research on large-scale 
volumetric simulations often employs data structures that reduce memory usage and accelerate neighbourhood access, 
such as Sparse Voxel Octrees~\cite{Laine2010} or the OpenVDB framework~\cite{Museth2013}. These approaches show 
how using locality-preserving representations can make simulations more scalable and offer guidance for extending neural 
CA to larger or more flexible volumetric setups.

\subsection{Training and Optimisation}

Because neural CA are differentiable, their parameters can be optimised using gradient-based methods to minimise a 
loss between the generated pattern and a target configuration~\cite{Mordvintsev2020}. Training generally involves 
unrolling the CA dynamics over multiple steps and applying backpropagation through time.

Typical loss functions include mean squared error (MSE) for continuous targets and intersection-over-union or other 
domain-specific metrics for structural objectives~\cite{Mordvintsev2020, Niklasson2021}. The choice of training regime—such 
as growth, persistence, or regeneration—affects the emergent behaviour and stability of the model. Optimisation often 
employs the Adam algorithm with learning rates around 0.001, sometimes supplemented by early stopping, data pooling, 
or logging frameworks~\cite{Sudhakaran2021}.

These training strategies enable the discovery of update rules capable of producing robust self-organising structures 
in both two- and three-dimensional systems~\cite{Rossi2021,Sudhakaran2021,Aadityaza2024}.

%=========================================================================
\chapter{Related Work}

The development of three-dimensional neural cellular automata draws on a lineage of contributions spanning mathematical biology,
continuous self-organising systems, and modern deep learning. This chapter surveys the most relevant prior work, organised thematically,
and concludes each section by relating the discussed contributions to the present thesis.

Some of the concepts discussed in this chapter were already introduced in Chapter~\ref{chap:theoretical-foundations}. There, they were 
presented to establish the formal and conceptual background required for the definition of three-dimensional neural cellular automata. 
In contrast, the present chapter considers the same works from the perspective of prior research contributions, experimental results, and 
architectural design choices that serve as a direct basis for the approach developed in this thesis. 


\chapter{Implementation}
TODO

\chapter{Testing}
TODO

\chapter{Experiments}
TODO

\chapter{Conclusion}
\label{zaver}

TODO

%=========================================================================

% For compilation piecewise (see projekt.tex), it is necessary to uncomment it
% \end{document}